%============================ DISEÑO (ENTREGA 2) ============================
\chapter{Diseño detallado}
\label{cap:diseno}

\section{Diagrama de componentes}
La Figura~\ref{fig:componentes} descompone el sistema en sus componentes de software y sus dependencias. En el servidor, \code{server.js} orquesta cuatro módulos especializados; en el firmware, el bucle principal coordina el escaneo, el portal cautivo y el enlace WebSocket.

\begin{figure}[H]
\centering
\begin{tikzpicture}[
  comp/.style={rectangle,draw,thick,rounded corners=2pt,align=center,font=\small,
               minimum width=2.7cm,minimum height=0.9cm},
  grp/.style={rectangle,draw,dashed,rounded corners=4pt,inner sep=8pt},
  arr/.style={-{Latex[length=2mm]},thick},
  lbl/.style={font=\scriptsize\itshape}
]
% Servidor
\node[comp,fill=eiloremerald!12] (srv) {server.js};
\node[comp,fill=eiloremerald!6,below left=1.0cm and 0.1cm of srv] (an) {lib/analysis};
\node[comp,fill=eiloremerald!6,below=1.0cm of srv] (ou) {lib/oui};
\node[comp,fill=eiloremerald!6,below right=1.0cm and 0.1cm of srv] (au) {lib/auth};
\node[comp,fill=eiloremerald!6,below=2.4cm of srv] (st) {lib/store};
\node[comp,fill=blue!7,below=0.7cm of st] (data) {data/*.json};

\draw[arr] (srv) -- (an); \draw[arr] (srv) -- (ou);
\draw[arr] (srv) -- (au); \draw[arr] (srv) -- (st);
\draw[arr] (st) -- (data);

% Firmware
\node[comp,fill=eilorcyan!12,right=4.2cm of srv] (loop) {loop()\\firmware};
\node[comp,fill=eilorcyan!6,above right=0.8cm and 0.2cm of loop] (scan) {Escaneo\\Wi-Fi/BLE};
\node[comp,fill=eilorcyan!6,right=1.0cm of loop] (portal) {Portal cautivo\\DNS+HTTP};
\node[comp,fill=eilorcyan!6,below right=0.8cm and 0.2cm of loop] (wsc) {Cliente\\WebSocket};
\draw[arr] (loop) -- (scan); \draw[arr] (loop) -- (portal); \draw[arr] (loop) -- (wsc);

% Enlace servidor-firmware
\draw[{Latex[length=2mm]}-{Latex[length=2mm]},thick,dashed] (srv) -- node[lbl,above]{WSS} (loop);
\end{tikzpicture}
\caption{Diagrama de componentes de software (servidor y firmware).}
\label{fig:componentes}
\end{figure}

\section{Flujo de datos: ciclo de escaneo}
La Figura~\ref{fig:seq-scan} muestra la secuencia de un ciclo de escaneo Wi\nobreakdash-Fi, desde la captura en el ESP32 hasta la actualización del dashboard.

\begin{figure}[H]
\centering
\begin{tikzpicture}[font=\small,
  msg/.style={-{Latex[length=2mm]},thick},
  self/.style={-{Latex[length=1.6mm]},thick},
  ml/.style={font=\scriptsize,fill=white,inner sep=1.5pt},
  act/.style={font=\scriptsize,align=center,fill=eiloremerald!8,draw=eiloremerald!40,rounded corners=2pt,inner sep=3pt}]
% Lifelines
\def\lx{0} \def\lserv{5.4} \def\ldash{10.8}
\node[font=\bfseries] (e) at (\lx,0) {ESP32};
\node[font=\bfseries] (s) at (\lserv,0) {Servidor};
\node[font=\bfseries] (d) at (\ldash,0) {Dashboard};
\foreach \x in {\lx,\lserv,\ldash} \draw[dashed,gray] (\x,-0.4) -- (\x,-7.0);
% 1. Escaneo local en el ESP32 (auto-mensaje)
\draw[self] (\lx,-0.9) -- ++(0.9,0) -- ++(0,-0.5) -- node[ml,right=2pt]{scanNetworks() · 13 canales} ++(-0.9,0);
% 2. ESP32 -> Servidor
\draw[msg] (\lx,-2.0) -- node[ml,above]{\{type:spectrum\_wifi, devices, channels\}} (\lserv,-2.0);
% 3. Procesamiento en el servidor
\node[act] at (\lserv,-2.8) {processIncomingScan()\\annotate · recalculateHeatmap\\analyzeChannels · detectEvilTwins};
% 4. Servidor -> Dashboard
\draw[msg] (\lserv,-3.9) -- node[ml,above]{event: scan\_update (heatmap, batch)} (\ldash,-3.9);
% 5. Render en el cliente
\node[act] at (\ldash,-4.7) {renderHeatmap()\\renderDevicesTable()};
% 6. Comando de control Dashboard -> Servidor -> ESP32
\draw[msg,eiloremerald!70!black] (\ldash,-5.7) -- node[ml,above]{set\_mode / create\_ap (+ token)} (\lserv,-5.7);
\draw[msg,eiloremerald!70!black] (\lserv,-6.4) -- node[ml,above]{comando reenviado por WSS} (\lx,-6.4);
\end{tikzpicture}
\caption{Diagrama de secuencia del ciclo de escaneo y control.}
\label{fig:seq-scan}
\end{figure}

\section{Flujo del portal cautivo con \emph{buffer} offline}
El registro de invitados sigue el flujo de la Figura~\ref{fig:seq-portal}. Si en el momento del envío no hay enlace con el servidor, el registro se almacena en una cola local y se reenvía automáticamente al reconectar, garantizando que no se pierdan altas.

\begin{figure}[H]
\centering
\begin{tikzpicture}[node distance=0.7cm,font=\small,
  st/.style={rectangle,rounded corners=2pt,draw,thick,align=center,minimum width=3.2cm,minimum height=0.8cm},
  dec/.style={diamond,aspect=2,draw,thick,align=center,font=\scriptsize,inner sep=1pt},
  arr/.style={-{Latex[length=2mm]},thick}]
\node[st,fill=eilorcyan!8] (a) {Cliente completa formulario};
\node[st,fill=eilorcyan!8,below=of a] (b) {ESP32 recibe POST /register};
\node[dec,below=of b,fill=yellow!12] (c) {¿Enlace WS?};
\node[st,fill=eiloremerald!10,below left=0.9cm and -0.6cm of c] (d) {Envía \{type:registration\}};
\node[st,fill=orange!12,below right=0.9cm and -0.6cm of c] (e) {Encola (máx. 20)};
\node[st,fill=eiloremerald!10,below=2.4cm of c] (f) {Servidor persiste y difunde};
\node[st,fill=orange!12,right=1.6cm of e] (g) {Reconexión $\Rightarrow$ flush};
\draw[arr] (a)--(b); \draw[arr] (b)--(c);
\draw[arr] (c) -| node[near start,above,font=\scriptsize]{sí} (d);
\draw[arr] (c) -| node[near start,above,font=\scriptsize]{no} (e);
\draw[arr] (d) |- (f);
\draw[arr] (e) -- (g); \draw[arr] (g) |- (f);
\end{tikzpicture}
\caption{Flujo de registro del portal cautivo con reenvío ante caída de enlace.}
\label{fig:seq-portal}
\end{figure}

\section{Modelo de datos}
El sistema maneja cuatro entidades principales, resumidas en la Tabla~\ref{tab:modelo}.

\begin{table}[H]
\centering
\caption{Entidades y atributos del modelo de datos}
\label{tab:modelo}
\footnotesize
\begin{tabularx}{\linewidth}{@{}l X@{}}
\toprule
\textbf{Entidad} & \textbf{Atributos principales} \\
\midrule
Dispositivo & \code{mac}, \code{type} (wifi/ble), \code{ssid}, \code{rssi}, \code{channel}, \code{encryption}, \code{distanceEst}, \code{vendor}, \code{randomMac}, \code{evilTwin}, \code{lastSeen} \\
Mapa de calor & \code{wifi[1..13]}: \{\code{freq}, \code{count}, \code{maxRssi}, \code{avgRssi}, \code{ssids}\}; \code{ble[37..39]}: \{\code{freq}, \code{count}, \code{maxRssi}\} \\
Muestra temporal & \code{t}, \code{wifi}, \code{ble}, \code{active}, \code{packets}, \code{rate} \\
Registro de invitado & \code{id}, \code{nombre}, \code{apellido}, \code{telefono}, \code{apSsid}, \code{source}, \code{ts} \\
\bottomrule
\end{tabularx}
\end{table}

\section{Especificación de interfaces}
\subsection{API REST}
La Tabla~\ref{tab:rest} documenta los catorce \emph{endpoints}, indicando el control de acceso aplicado (sesión de usuario, token de dispositivo o público).

\begin{table}[H]
\centering
\caption{Endpoints de la API REST}
\label{tab:rest}
\footnotesize
\begin{tabularx}{\linewidth}{@{}l l X l@{}}
\toprule
\textbf{Método} & \textbf{Ruta} & \textbf{Función} & \textbf{Acceso} \\
\midrule
POST & /api/login & Iniciar sesión, emite token & Público \\
POST & /api/logout & Cerrar sesión & Sesión \\
GET  & /api/session & Validar token & Público \\
POST & /api/scan & Ingesta de escaneos & Dispositivo \\
GET  & /api/scan & Estado y dispositivos & Público \\
POST & /api/mode & Cambiar modo del ESP32 & Sesión \\
POST & /api/ap & Desplegar SoftAP & Sesión \\
GET  & /api/heatmap & Mapa de calor de canales & Público \\
GET  & /api/timeseries & Serie temporal & Público \\
POST & /api/registration & Alta de invitado & Dispositivo \\
GET  & /api/registrations & Listar registros & Sesión \\
GET  & /api/registrations.csv & Exportar CSV & Sesión \\
POST & /api/registrations/clear & Borrar registros & Sesión \\
POST & /api/clear & Reiniciar datos & Sesión \\
\bottomrule
\end{tabularx}
\end{table}

\subsection{Mensajería WebSocket}
El canal WebSocket transporta mensajes de cliente a servidor ---\code{get\_state}, \code{register}/\code{heartbeat}, escaneos (\code{type} wifi/ble), \code{set\_mode}/\code{toggle\_mode}, \code{create\_ap}, \code{registration} y \code{ping}--- y eventos de servidor a cliente ---\code{initial\_state}, \code{scan\_update}, \code{esp32\_status\_change}, \code{telemetry\_log}, \code{timeseries\_sample}, \code{registration}, \code{registrations\_cleared}, \code{clear}, \code{auth\_error} y \code{pong}---. Las acciones de control validan el token de sesión y la ingesta valida el token de dispositivo antes de ejecutarse.
